\subsection{Belief revision, argumentation and the status of negative evidence}

The belief-maintenance lineage explains why revision cannot be implemented as destructive overwrite. Doyle's truth-maintenance system records justifications so downstream beliefs can be withdrawn when their supports fail, while assumption-based TMS preserves multiple environments rather than forcing a single consistent context \cite{doyle1979,dekleer1986}. AGM revision supplies a complementary axiomatic view of minimal change under incorporation of new information \cite{agm1985}. Abstract argumentation adds a different object: a network in which arguments can attack one another and acceptability depends on the chosen semantics \cite{dung1995}. These frameworks are not interchangeable. TMS emphasizes dependency maintenance, AGM emphasizes rational belief-set change, and argumentation emphasizes dialectical support/attack structure. \RAKL{} borrows the need to preserve dependencies and incompatibilities, but its canonical objects are scientific claims scoped by evidence, context and authority rather than unqualified propositions.

This difference is most visible for negative evidence. A null, refutation or failed mapping has two temporal identities: it is evidence about the historical state of inquiry, and it may or may not remain decisive for a later, differently scoped claim. Deleting the old object when a successor arrives destroys information needed to audit cherry-picking, repeated testing, route resurrection and model-family evolution. The rule
\[
\mathcal H^-_t\subseteq \mathcal H^-_{t+1}
\]
therefore means preservation of the event and its original scope, not that every old negative conclusion remains eternally applicable. A later observation can narrow, supersede or contextualize the interpretation while leaving the receipt addressable.

This is also a response to a known scientific-selection problem. The ``file drawer'' literature shows why invisibility of null results can distort the apparent evidence landscape \cite{rosenthal1979,franco2014}. In \RAKL{}, negative history is not merely a publication-policy preference. It is a state variable used by search, saturation and self-evolution. If a candidate mechanism has already failed under a particular regime, a later generator must either explain why the new context is materially different or inherit that negative evidence. If a method repair was tried and failed, the next Self-\RAKL{} generation cannot erase the failure to make its lineage look monotone.

Model pluralism further motivates resisting forced synthesis. Veit argues that many scientific purposes require sets of models whose functions depend on context and history rather than one canonical model \cite{veit2020pluralism}. \RAKL{} gives that pluralism an operational destination: a plural atlas can be a successful outcome when different local models remain useful under distinct contexts, while an identified set is appropriate when the models compete on the same target but observations cannot distinguish them. The two should not be confused. Plurality can arise from legitimate perspectival specialization; non-identification arises from insufficient discriminatory evidence.

\subsection{Scientific revision as an auditable transaction}

The preceding traditions suggest a stronger implementation rule: scientific revision should be represented as a transaction with explicit preconditions, postconditions and preserved supersession lineage. A conventional belief store can replace proposition \(p\) with \(\neg p\). A scientific atlas must additionally answer which population, measurement, assumptions and evidence made each version applicable. The transaction object is therefore richer than a truth-value flip. One useful abstract form is
\[
\rho=(c_{\mathrm{old}},c_{\mathrm{new}},e,\gamma,\tau,\omega),
\]
where \(e\) is the evidence that triggered revision, \(\gamma\) its context, \(\tau\) the typed relation between old and new claims, and \(\omega\) the authority effect that the evidence is licensed to create. The old claim remains historically addressable even when its current applicability becomes empty.

This representation distinguishes at least five revision patterns. \emph{Refutation} records evidence incompatible with a claim under aligned context. \emph{Scope restriction} preserves the claim on a narrower domain after a counterexample appears elsewhere. \emph{Supersession} replaces an effective or approximate law with a successor while preserving the regime in which the old law remains adequate. \emph{Decomposition} splits one coarse claim into several context-specific claims after hidden heterogeneity is discovered. \emph{Authority downgrade} retains the descriptive content but withdraws a stronger mechanism, identification or decision certificate when its prerequisite fails. These outcomes are scientifically different even if a final prose summary would say simply that ``the model was updated.''

Dependency propagation is also typed. If claim \(b\) was derived from \(a\), refuting \(a\) may invalidate \(b\). If \(b\) was independently measured and merely consistent with \(a\), the same refutation need not remove \(b\). A provenance graph that records only citation ancestry cannot always make this distinction; the transition relation must say whether an edge is justificatory, derivational, contextual, analogical or merely navigational. This is where truth-maintenance ideas meet epistemic provenance: dependencies matter because they determine which descendants should be reconsidered when evidence changes \cite{doyle1979,dekleer1986}.

Argumentation adds another useful boundary. An attack edge can express that two claims cannot both retain a particular authority under a registered context, but an attack is not itself a refutation. A methodological critique may attack an inference while leaving the measurement intact. A new experiment may attack one mechanism but leave a predictive relation untouched. A conflict between two observational studies may be suspended until population and measurement alignment is established. The atlas therefore does not ask only ``which argument wins?''; it asks which coordinates of which scientific objects remain licensed after the conflict is typed \cite{dung1995}.

Revision has a chronology obligation. The system must know whether a hypothesis, threshold, covariate set or interpretation was fixed before the evidence that appears to support it. Otherwise the same data can play both proposal and confirmation roles without disclosure. This is why negative history, preregistration and exact candidate identity belong to the same control architecture. They are different ways of preserving the temporal direction of scientific justification.

The transaction view also explains why a paper should not silently rewrite its own past. If a manuscript once claimed that the generic atom--witness structure was a mathematical lattice and later analysis showed that only a compatibility complex had been established, the correct repair is not merely a global search-and-replace. The project should preserve the terminology error as method history, document the narrower structure, and state which downstream claims changed. The current compatibility-complex repair follows that pattern. Scientific self-correction is more informative when the reader can see what was corrected and why.

Finally, auditable revision provides a natural interface to sensitivity analysis. Instead of asking whether one final conclusion is ``robust'', the system can replay the transition under alternative registered assumptions or context mappings and observe which authority coordinates change. A conclusion that survives many plausible revisions has a different scientific profile from one that depends on a single fragile alignment, even if both currently have the same point estimate. Orion treats that profile as structured evidence rather than compressing it into a confidence adjective.

\subsection{Why evidential selection pressure belongs inside the model of research}

Scientific literature is not an unbiased sample of all completed inquiries. Selective publication, flexible analysis and repeated exploration can change the distribution of claims that enter a search system \cite{rosenthal1979,franco2014,ioannidis2005}. A literature agent that retrieves published claims perfectly can therefore inherit upstream selection bias perfectly. Evidence governance begins before language-model reasoning.

\RAKL{} does not attempt a universal correction for publication bias. Instead it records relevant selection information when available and avoids treating publication count as independent confirmation. Pre-registration status, registered reports, availability of null results, dataset ancestry, replication status and analysis chronology can all alter how a claim should be interpreted. These become context/evidence attributes rather than one global quality label.

The same reasoning applies to benchmark ecosystems. Scientific-agent benchmarks are published objects selected by designers, and method developers adapt to visible tasks. A result on a saturated public benchmark can overstate transfer if the benchmark has become part of the development environment. Fresh-assurance tasks and frozen evaluation chronology are intended to protect against this adaptive exposure at the method level.

This is another reason negative history has value beyond transparency. Failed or null results change the selection model of what has been tried. Preserving them can prevent repeated rediscovery of a dead end and can reveal that an apparently dominant positive literature sits on top of many unreported or locally stored failures. In a research organization, the internal atlas can therefore contain scientifically valuable evidence that never became a conventional publication.

The authority system treats these factors as modifiers and blockers rather than as automatic verdicts. A preregistered study can be badly measured; a post hoc analysis can reveal a real phenomenon; a highly replicated result can still fail to transport. The purpose is to preserve the variables needed for scientific judgment and downstream sensitivity analysis.

\subsection{Provenance and machine-readable research objects}
Database provenance work had already made the origin and derivation of view data a first-class problem, including the distinction between ``why'' and ``where'' provenance \cite{buneman2001}. The W3C PROV family formalizes entities, activities, agents and derivations as provenance objects \cite{moreau2013}; FAIR articulates findability, accessibility, interoperability and reuse as stewardship goals \cite{wilkinson2016}; RO-Crate packages data, software and methods with machine-readable metadata and relationships \cite{soilandreyes2022}; and nanopublications attach provenance to granular scientific assertions \cite{kuhn2018}. Agent-workflow provenance and verifiable agent-native publishing are also active contemporary directions \cite{souza2025,dogah2026}. Orion therefore does not claim generic provenance or atomic statement packaging. Its narrower addition is to bind provenance to context-scoped scientific authority, negative history, evidence-lineage-aware saturation and target-conditioned context compilation. Computational reproducibility is likewise an established discipline: reproducible-computation rules emphasize recording how each result was produced \cite{sandve2013}, while good-enough scientific-computing practice emphasizes data, software, collaboration and project organization \cite{wilson2017}. Orion's receipts, environment attestation and artifact manifests are an implementation of this broader reproducibility requirement, not a novelty claim by themselves.

Recent agent-verification work also shows why pooled factual support is insufficient. ProvenanceGuard evaluates whether an atomic claim is supported by the particular source to which it is attributed, exposing cross-source conflation as a distinct error from ordinary factuality \cite{alvarez2026provenanceguard}. A broader survey of evidence tracing and execution provenance likewise separates retrieved evidence, tool outputs, memory, intermediate claims, actions and final answers rather than treating an agent trace as one undifferentiated rationale \cite{wang2026traces}. Orion adopts the same granularity but adds a scientific-authority boundary: a correctly attributed source can still be weak, context-incompatible, derivative of the same evidence root, or irrelevant to a mechanism claim. Source ownership is therefore necessary for some updates but never sufficient for arbitrary authority escalation.

\paragraph{Provenance standards and the novelty boundary.}
Provenance representation itself is established infrastructure. W3C PROV-O, for example, supplies a general ontology for interchanging provenance across heterogeneous systems \cite{w3cprov2013}. \RAKL{} therefore does not claim novelty for provenance graphs or source ancestry. Its narrower contribution is to couple source-rehydratable provenance to context-scoped scientific objects, typed relation witnesses, multi-axis authority and gated canonical updates; provenance visibility alone neither establishes truth nor licenses a mechanism or identification claim.
